\section{Hồ sơ tái lập và truy nguyên}
\label{sec:reproducibility}

\subsection{Nguồn bằng chứng}

Các kết quả được tổng hợp từ bốn nhóm bằng chứng: dữ liệu và phép chia đã kiểm định; artifact huấn luyện/dự đoán; báo cáo thống kê bắt cặp; và gói công bố gồm bảng, hình cùng ma trận bằng chứng. Vai trò của truy nguyên là bảo đảm mỗi con số và kết luận có thể quay lại đúng artifact, không phải làm tăng điểm số hay thay thế xác nhận trên mẫu độc lập.

\begin{table}[H]
\centering
\caption{Các nhóm bằng chứng được sử dụng trong báo cáo}
\label{tab:reproducibility_evidence}
\begin{tabularx}{\textwidth}{lY}
\toprule
\textbf{Nhóm} & \textbf{Vai trò} \\
\midrule
Dữ liệu và phép chia & Xác nhận đối tượng, bản ghi, nhãn và nguyên tắc tách mẫu \\
Kết quả huấn luyện & Xác nhận mô hình đã chọn, dự đoán, chỉ số và sự tương thích giữa cấu hình \\
Phân tích thống kê & Cung cấp chênh lệch, khoảng tin cậy và kiểm định bắt cặp \\
Gói công bố & Cung cấp bảng, hình, ma trận bằng chứng và kiểm tra toàn vẹn độc lập \\
\bottomrule
\end{tabularx}
\end{table}

\subsection{Nguyên tắc kiểm tra}

Kiểm tra tái lập được thực hiện theo ba lớp: (1) cấu trúc, độ phủ và khóa bắt cặp; (2) phiên bản, mã băm và tính đầy đủ của artifact; (3) sự khớp nhau giữa bảng/hình/tuyên bố và dữ liệu nguồn. Các phép kiểm tra chỉ đọc artifact đã khóa, không huấn luyện lại mô hình hoặc mở lại test để chọn kết quả.

Lần lặp seed 123 dùng cùng dữ liệu, phép chia và họ cấu hình. Các seed được báo cáo riêng, không gộp $p$-value và không được xem như mẫu đầy đủ của biến thiên khởi tạo. Các phân tích SHHS thành phần, E3--E2 và E4 được ghi rõ là thứ cấp hoặc hậu nghiệm trên cohort đã mở.

\subsection{Tái phân tích E6 và tình trạng provenance SHHS1}

Phân tích E6 theo lớp được tính trực tiếp từ 180 artifact dự đoán chuyển miền không cập nhật trọng số đã khóa trên ổ ngoài. Chiến dịch được đánh dấu hoàn tất, test gate đạt, toàn bộ 180 mã băm artifact E6 khớp run manifest, và ma trận nhầm lẫn tính lại theo từng bản ghi khớp giá trị trong manifest. Vì vậy, các chỉ số E6 mới trong báo cáo là phép tái phân tích mô tả của dự đoán đã tồn tại, không phải một lượt suy luận mới.

Theo hồ sơ kiểm toán đã lưu, dự đoán E6 và chỉ số tính lại đã được xác minh ở cấp artifact. Trong lượt rà soát hiện tại, run manifest gốc trên ổ ngoài vẫn truy cập được và có SHA-256 \seqsplit{f9cd5ebbd20f26b188b5dc13ac6e417ff8ef0fa8dcae78760cfcb27940bf58cf}. Protocol hash ghi trong manifest (\seqsplit{165d7cdf614ff071da7bd5ca94eb4e52dd8bee1ce5eafb712c2c8a0d0550fe93}) khớp snapshot lịch sử \texttt{configs/shhs\_zero\_shot\_v1.json}. Tệp \texttt{configs/shhs\_v1\_protocol.json} hiện có hash \seqsplit{9541e2334cdae98b5d36b95a2656993cdaaffb35a3160dad70af11d14b653fe9} vì là hồ sơ mở rộng sau chạy; không thay hash lịch sử bằng hash này. Khớp hash xác nhận danh tính tệp và liên kết artifact, không tự chứng minh đăng ký trước công khai, thời điểm khóa độc lập hoặc khả năng tái sinh toàn bộ chiến dịch. Lượt chỉnh sửa báo cáo không huấn luyện hay tạo prediction mới.

\subsection{Tình trạng dữ liệu dẫn xuất}

Snapshot dữ liệu dẫn xuất gồm 765 tệp thuộc năm biến thể tiền xử lý. Kiểm toán độc lập xác nhận 765/765 tệp khớp SHA-256 trong manifest nội dung, có metadata ZIP chuẩn hóa, đọc được và nhất quán về nhãn, vị trí thời gian, hình dạng cùng quan hệ biến đổi khoa học. Mã băm container lịch sử được bảo toàn trong trường truy nguyên riêng.

Khóa byte của snapshot bảo đảm tính toàn vẹn khi lưu trữ và trao đổi, nhưng không thay thế phép tái sinh độc lập từ EDF gốc trong môi trường phần mềm đã khóa. Gói mã nguồn đi kèm giao thức và công cụ kiểm định cho bước tái sinh đó.

\subsection{Biên diễn giải khoa học}

\begin{table}[H]
\centering
\caption{Phạm vi diễn giải của các kết quả chính}
\label{tab:claim_boundaries}
\begin{tabularx}{\textwidth}{YY}
\toprule
\textbf{Kết luận được hỗ trợ} & \textbf{Diễn giải vượt quá bằng chứng} \\
\midrule
Các phép thay cấu hình E1/E2 chưa cho ưu thế ổn định & Đã cô lập kiến trúc, hoặc ResNet-1D/TCN vượt trội phổ quát \\
ResNet-1D--TCN nhanh hơn và gọn hơn về vận hành & Mô hình nhẹ hơn về tham số hoặc bộ nhớ \\
E3 cao hơn E0/E6 trên mẫu SHHS1 đã khóa & E3 đã giải quyết dịch chuyển miền hoặc vượt SOTA \\
E3--E2/E4 cho lợi ích của cấu hình tiền xử lý trên mẫu đã khảo sát & Một thao tác riêng là nguyên nhân hoặc trục tiền xử lý luôn ưu việt \\
N3$\rightarrow$N2 tái diễn ở E0/E3 và có đòn bẩy phản thực lớn nhất trên E3 & E3 gây ra lỗi này hoặc sửa kênh này chắc chắn đạt \mf\ 0,7444 \\
E6 không sửa được lỗi N3 trong pipeline này & Biên độ tuyệt đối không liên quan đến N3 \\
Gate 8 chưa cho thấy giá trị biên của P/N & Ngữ cảnh thời gian hoàn toàn vô dụng \\
Silhouette không phù hợp để chọn encoder ở đây & Embedding CNN tốt hơn ResNet cho mọi nhiệm vụ \\
Các điều kiện chưa được kiểm định tương đương & Giá trị $p$ lớn chứng minh tương đương \\
\bottomrule
\end{tabularx}
\end{table}

\subsection{Thông tin tác giả và các xác nhận trước công bố}

Ngô Nhật Quân là tác giả đứng đầu và tác giả liên hệ, email \texttt{23521258@gm.uit.edu.vn}; Phạm Thái Sơn là đồng tác giả thứ hai. Hai tác giả đã đọc và duyệt bản cuối. Ngô Nhật Quân: Conceptualization, Methodology, Software, Data curation, Formal analysis, Investigation, Visualization, Writing -- original draft, Writing -- review and editing. Phạm Thái Sơn: Validation, Investigation, Visualization, Writing -- review and editing.

Các tác giả trân trọng cảm ơn Nguyễn Hồ Duy Trí đã hướng dẫn học thuật, định hướng đề tài nghiên cứu và góp ý xây dựng cho các phiên bản trước của bản thảo.

Tác giả xác nhận đã xin được quyền truy cập SHHS qua NSRR và xác nhận phân tích thứ cấp dữ liệu đã khử định danh không cần xét duyệt đạo đức tại đơn vị áp dụng. Nghiên cứu nhận khoảng 6.000.000 VND hỗ trợ sinh viên từ Trường Đại học Công nghệ Thông tin, Đại học Quốc gia Thành phố Hồ Chí Minh; hai tác giả tuyên bố không có xung đột lợi ích. Repository mã nguồn là \url{https://github.com/quanuiter/SleepTCN}; báo cáo không công bố dữ liệu hạn chế hoặc mã định danh SHHS.

OpenAI Codex chỉ hỗ trợ biên tập, tổ chức nội dung, kiểm tra nhất quán với code/artifact và dựng hình từ các tổng hợp số liệu đã lưu; không huấn luyện thêm mô hình hoặc tạo dữ liệu thực nghiệm. Tác giả không dùng công cụ AI nào khác, hai tác giả đã đọc và duyệt bản cuối và chịu trách nhiệm về nội dung trước công bố.
